\documentclass[10pt,a4paper]{amsart}
\usepackage{amsmath,amssymb,amsthm,graphicx}
% Retained minimal preamble for the Stacks Project chapter more-algebra.tex (More on Algebra).
% The chapter file pinned-original-source.tex carries only \input{preamble} on its first line, so no
% preamble text lives inside the chapter itself.  The Stacks Project uses one single shared preamble for
% every chapter; the snapshot stages it identically as sources/stacks-chow/preamble.tex and
% sources/stacks-etale-cohomology/preamble.tex.  The block below is copied mechanically, line for line,
% from sources/stacks-chow/preamble.tex lines 14-18 (the Xy-pic setup) and 154-217 (the native theorem
% declarations together with the mathematical macros).  Nothing was added, removed or redefined; the
% wrapper class and the amsmath/amssymb/amsthm/graphicx base packages are supplied by the assembler.
% Chapter: more-algebra.tex -- More on Algebra.
% For commutative diagrams we use Xy-pic
\usepackage[all]{xy}

% We use 2cell for 2-commutative diagrams.
\xyoption{2cell}

% Theorem environments.
%
\theoremstyle{plain}
\newtheorem{theorem}[subsection]{Theorem}
\newtheorem{proposition}[subsection]{Proposition}
\newtheorem{lemma}[subsection]{Lemma}

\theoremstyle{definition}
\newtheorem{definition}[subsection]{Definition}
\newtheorem{example}[subsection]{Example}
\newtheorem{exercise}[subsection]{Exercise}
\newtheorem{situation}[subsection]{Situation}

\theoremstyle{remark}
\newtheorem{remark}[subsection]{Remark}
\newtheorem{remarks}[subsection]{Remarks}

\numberwithin{equation}{subsection}

% Macros
%
\def\lim{\mathop{\mathrm{lim}}\nolimits}
\def\colim{\mathop{\mathrm{colim}}\nolimits}
\def\Spec{\mathop{\mathrm{Spec}}}
\def\Hom{\mathop{\mathrm{Hom}}\nolimits}
\def\Ext{\mathop{\mathrm{Ext}}\nolimits}
\def\SheafHom{\mathop{\mathcal{H}\!\mathit{om}}\nolimits}
\def\SheafExt{\mathop{\mathcal{E}\!\mathit{xt}}\nolimits}
\def\Sch{\mathit{Sch}}
\def\Mor{\mathop{\mathrm{Mor}}\nolimits}
\def\Ob{\mathop{\mathrm{Ob}}\nolimits}
\def\Sh{\mathop{\mathit{Sh}}\nolimits}
\def\NL{\mathop{N\!L}\nolimits}
\def\CH{\mathop{\mathrm{CH}}\nolimits}
\def\proetale{{pro\text{-}\acute{e}tale}}
\def\etale{{\acute{e}tale}}
\def\QCoh{\mathit{QCoh}}
\def\Ker{\mathop{\mathrm{Ker}}}
\def\Im{\mathop{\mathrm{Im}}}
\def\Coker{\mathop{\mathrm{Coker}}}
\def\Coim{\mathop{\mathrm{Coim}}}

% Boxtimes
%
\DeclareMathSymbol{\boxtimes}{\mathbin}{AMSa}{"02}

%
% Macros for moduli stacks/spaces
%
\def\QCohstack{\mathcal{QC}\!\mathit{oh}}
\def\Cohstack{\mathcal{C}\!\mathit{oh}}
\def\Spacesstack{\mathcal{S}\!\mathit{paces}}
\def\Quotfunctor{\mathrm{Quot}}
\def\Hilbfunctor{\mathrm{Hilb}}
\def\Curvesstack{\mathcal{C}\!\mathit{urves}}
\def\Polarizedstack{\mathcal{P}\!\mathit{olarized}}
\def\Complexesstack{\mathcal{C}\!\mathit{omplexes}}
% \Pic is the operator that assigns to X its picard group, usage \Pic(X)
% \Picardstack_{X/B} denotes the Picard stack of X over B
% \Picardfunctor_{X/B} denotes the Picard functor of X over B
\def\Pic{\mathop{\mathrm{Pic}}\nolimits}
\def\Picardstack{\mathcal{P}\!\mathit{ic}}
\def\Picardfunctor{\mathrm{Pic}}
\def\Deformationcategory{\mathcal{D}\!\mathit{ef}}
\usepackage{cleveref}
\crefname{originallemma}{original lemma}{original lemmas}
\Crefname{originallemma}{Original lemma}{Original lemmas}
\makeatletter
\newlabel{more-algebra-lemma-relative-finitely-presented-module-over-fibre-product}{{08KP}{source}{}{sourcepointer-1}{}}
\newlabel{more-algebra-lemma-relative-finitely-presented-module-over-fibre-product@cref}{{[originallemma][0][]08KP}{[1][0][]source}{}{}{}}
\makeatother
\crefname{originallemma}{original lemma}{original lemmas}
\Crefname{originallemma}{Original lemma}{Original lemmas}
\makeatletter
\newlabel{lemma-relative-module-over-fibre-product}{{08KM}{source}{}{sourcepointer-2}{}}
\newlabel{lemma-relative-module-over-fibre-product@cref}{{[originallemma][0][]08KM}{[1][0][]source}{}{}{}}
\makeatother
\crefname{originallemma}{original lemma}{original lemmas}
\Crefname{originallemma}{Original lemma}{Original lemmas}
\makeatletter
\newlabel{algebra-lemma-finite-finitely-presented-extension}{{0564}{source}{}{sourcepointer-3}{}}
\newlabel{algebra-lemma-finite-finitely-presented-extension@cref}{{[originallemma][0][]0564}{[1][0][]source}{}{}{}}
\makeatother
\crefname{originallemma}{original lemma}{original lemmas}
\Crefname{originallemma}{Original lemma}{Original lemmas}
\makeatletter
\newlabel{lemma-relative-finite-module-over-fibre-product}{{08IH}{source}{}{sourcepointer-4}{}}
\newlabel{lemma-relative-finite-module-over-fibre-product@cref}{{[originallemma][0][]08IH}{[1][0][]source}{}{}{}}
\makeatother
\crefname{originallemma}{original lemma}{original lemmas}
\Crefname{originallemma}{Original lemma}{Original lemmas}
\makeatletter
\newlabel{lemma-flat-module-over-fibre-product}{{0D2I}{source}{}{sourcepointer-5}{}}
\newlabel{lemma-flat-module-over-fibre-product@cref}{{[originallemma][0][]0D2I}{[1][0][]source}{}{}{}}
\makeatother
\crefname{originallemma}{original lemma}{original lemmas}
\Crefname{originallemma}{Original lemma}{Original lemmas}
\makeatletter
\newlabel{algebra-lemma-extension}{{0519}{source}{}{sourcepointer-6}{}}
\newlabel{algebra-lemma-extension@cref}{{[originallemma][0][]0519}{[1][0][]source}{}{}{}}
\makeatother
\crefname{originallemma}{original lemma}{original lemmas}
\Crefname{originallemma}{Original lemma}{Original lemmas}
\makeatletter
\newlabel{algebra-lemma-flat-tor-zero}{{00HL}{source}{}{sourcepointer-7}{}}
\newlabel{algebra-lemma-flat-tor-zero@cref}{{[originallemma][0][]00HL}{[1][0][]source}{}{}{}}
\makeatother
\title{A finitely presented module over a fibre product of rings}
\author{The Stacks Project Authors}
\date{stacks-project:more-algebra.tex:local-snapshot-20261012; GNU FDL 1.2 or later}
\begin{document}
\maketitle
\setcounter{subsection}{1}
\noindent\textit{Editorial note.} All delivered mathematical text is the author's own, copied from the pinned Stacks Project chapter file more-algebra.tex; no statement or proof step was rewritten, reordered or removed.\par\medskip
\noindent\textit{Editorial note.} The chapter file pinned-original-source.tex carries only \textbackslash{}input\{preamble\} on its first line, so no preamble text lives inside the chapter. The retained file sources/stacks-more-algebra/retained-preamble.tex is a mechanical, line-for-line copy of the project's single shared preamble as staged in the snapshot (Xy-pic setup plus theorem declarations and mathematical macros). Nothing was added, removed or redefined, and the excerpt serves only as preamble, never as a span coordinate.\par\medskip
\noindent\textit{Editorial note.} Both counted spans start at the \textbackslash{}begin line of their environment, so the statement and the author's proof are delivered whole; surrounding prose of the section is deliberately outside the counted spans and is carried as standard prerequisites instead.\par\medskip
\noindent\textit{Editorial note.} The citation fields (source\_id, source\_version, authors, source\_url, license.name, source\_title) are identical to those of every other card this lane has produced from this chapter, so the batch selector treats them as one source.\par\medskip
\noindent\textit{Allowed standard prerequisites (editorial).}\begin{itemize}
\item Fibre product of rings and modules over it, as set up at the start of the section. source lines 521-1097, before the counted statement. Supplies the rings and the description of modules over the fibre product.
\item Finite presentation of a module and the presentation matrix attached to it. used throughout the proof, lines 1111-1156. Provides the object that is decomposed into the two blocks.
\item Base change of a module along the two projections of the fibre product. used at proof lines 1111-1156. Supplies the two modules over the factors and the identification between them.
\end{itemize}\medskip
\ifdefined\hypertarget\hypertarget{sourcepointer-1}{}\fi
\noindent\textit{Source reference.} Original lemma 08KP: focused-seeds/AGP/sources/stacks-snapshot-20261011/tags.txt, source lines 10892--10892. The lemma's own canonical tag, quoted so that readers can resolve the label inside the delivered text; it points at the counted statement itself and carries no proof content.\par\medskip
\ifdefined\hypertarget\hypertarget{sourcepointer-2}{}\fi
\noindent\textit{Source reference.} Original lemma 08KM: focused-seeds/AGP/sources/stacks-snapshot-20261011/tags.txt, source lines 10890--10890. Navigation pointer to a statement the counted proof refers to but does not contain; nothing load-bearing is delegated to it.\par\medskip
\ifdefined\hypertarget\hypertarget{sourcepointer-3}{}\fi
\noindent\textit{Source reference.} Original lemma 0564: focused-seeds/AGP/sources/stacks-snapshot-20261011/tags.txt, source lines 6697--6697. Navigation pointer to a statement the counted proof refers to but does not contain; nothing load-bearing is delegated to it.\par\medskip
\ifdefined\hypertarget\hypertarget{sourcepointer-4}{}\fi
\noindent\textit{Source reference.} Original lemma 08IH: focused-seeds/AGP/sources/stacks-snapshot-20261011/tags.txt, source lines 10815--10815. Navigation pointer to a statement the counted proof refers to but does not contain; nothing load-bearing is delegated to it.\par\medskip
\ifdefined\hypertarget\hypertarget{sourcepointer-5}{}\fi
\noindent\textit{Source reference.} Original lemma 0D2I: focused-seeds/AGP/sources/stacks-snapshot-20261011/tags.txt, source lines 16383--16383. Navigation pointer to a statement the counted proof refers to but does not contain; nothing load-bearing is delegated to it.\par\medskip
\ifdefined\hypertarget\hypertarget{sourcepointer-6}{}\fi
\noindent\textit{Source reference.} Original lemma 0519: focused-seeds/AGP/sources/stacks-snapshot-20261011/tags.txt, source lines 6527--6527. Navigation pointer to a statement the counted proof refers to but does not contain; nothing load-bearing is delegated to it.\par\medskip
\ifdefined\hypertarget\hypertarget{sourcepointer-7}{}\fi
\noindent\textit{Source reference.} Original lemma 00HL: focused-seeds/AGP/sources/stacks-snapshot-20261011/tags.txt, source lines 647--647. Navigation pointer to a statement the counted proof refers to but does not contain; nothing load-bearing is delegated to it.\par\medskip
\begin{lemma}
\label{lemma-relative-finitely-presented-module-over-fibre-product}
Let $A, A', B, B', C, C', D, D', I, M', M, N, \varphi$ be as in
Lemma \ref{lemma-relative-module-over-fibre-product}. If
\begin{enumerate}
\item $N$ is finitely presented over $D$ and flat over $B$,
\item $M'$ finitely presented over $C'$ and flat over $A'$, and
\item the ring map $B' \to D'$ factors as $B' \to D'' \to D'$
with $B' \to D''$ flat and $D'' \to D'$ of finite presentation,
\end{enumerate}
then $N' = N \times_M M'$ is finitely presented over $D'$.
\end{lemma}

\begin{proof}
Choose a surjection $D''' = D''[x_1, \ldots, x_n] \to D'$ with
finitely generated kernel $J$.
By Algebra, Lemma \ref{algebra-lemma-finite-finitely-presented-extension}
it suffices to show that $N'$ is finitely presented as a
$D'''$-module. Moreover, $D''' \otimes_{B'} B \to D' \otimes_{B'} B = D$
and $D''' \otimes_{B'} A' \to D' \otimes_{B'} A' = C'$ are surjections
whose kernels are generated by the image of $J$, hence $N$ is a
finitely presented $D''' \otimes_{B'} B$-module and
$M'$ is a finitely presented $D''' \otimes_{B'} A'$-module by
Algebra, Lemma \ref{algebra-lemma-finite-finitely-presented-extension}
again. Thus we may replace $D'$ by $D'''$ and $D$ by
$D''' \otimes_{B'} B$, etc. Since $D'''$ is
flat over $B'$, it follows that we may assume that $B' \to D'$ is flat.

\medskip\noindent
Assume $B' \to D'$ is flat.
By Lemma \ref{lemma-relative-finite-module-over-fibre-product}
the module $N'$ is finite over $D'$. Choose a surjection
$(D')^{\oplus n} \to N'$ with kernel $K'$. By base change we obtain maps
$D^{\oplus n} \to N$, $(C')^{\oplus n} \to M'$, and $C^{\oplus n} \to M$
with kernels $K_D$, $K_{C'}$, and $K_C$. There is a canonical map
$$
K' \longrightarrow K_D \times_{K_C} K_{C'}
$$
On the other hand, since $N' = N \times_M M'$ and
$D' = D \times_C C'$ (by Lemma \ref{lemma-flat-module-over-fibre-product};
applied to the flat $B'$-module $D'$)
there is also a
canonical map $K_D \times_{K_C} K_{C'} \to K'$ inverse to the displayed
arrow. Hence the displayed map is an isomorphism. By
Algebra, Lemma \ref{algebra-lemma-extension}
the modules $K_D$ and $K_{C'}$ are finite. We conclude from
Lemma \ref{lemma-relative-finite-module-over-fibre-product}
that $K'$ is a finite $D'$-module provided that $K_D \to K_C$ and
$K_{C'} \to K_C$ induce isomorphisms
$K_D \otimes_B A = K_C = K_{C'} \otimes_{A'} A$.
This is true because the flatness assumptions implies the sequences
$$
0 \to K_D \to D^{\oplus n} \to N \to 0
\quad\text{and}\quad
0 \to K_{C'} \to (C')^{\oplus n} \to M' \to 0
$$
stay exact upon tensoring, see
Algebra, Lemma \ref{algebra-lemma-flat-tor-zero}.
\end{proof}

\end{document}
